\documentclass[11pt]{article}
% Shared design for the two minimal papers.
% It follows the layout of the guide "The Riemann Hypothesis and its generalizations"
% (KNT with Fable, 9 November 2026): a designed title block, a short story list at the start,
% coloured headings, boxed statements, idea boxes, and annotated figures.
%
% A paper sets \papertitle, \papersubtitle and \paperrunning before \begin{document}
% and calls \maketitleblock after it.

\usepackage[letterpaper,left=1.12in,right=1.12in,top=1.05in,bottom=1.1in]{geometry}
\usepackage{amsmath,mathtools,amssymb,amsthm}
\usepackage{newpxtext}
\usepackage[varbb]{newpxmath}
\usepackage{fontspec}
\setsansfont{texgyreheros}[Extension=.otf, UprightFont=*-regular, BoldFont=*-bold,
  ItalicFont=*-italic, BoldItalicFont=*-bolditalic, Scale=0.88]
\usepackage{microtype}
\usepackage{parskip}
\usepackage{xcolor}
\usepackage{booktabs,array}
\usepackage{enumitem}
\usepackage{multicol}
\usepackage[titles]{tocloft}
\usepackage{tikz}
\usetikzlibrary{decorations.markings,arrows.meta,positioning,calc}
\usepackage[most]{tcolorbox}
\usepackage{fancyhdr}
\usepackage{titlesec}
\usepackage[font=small,labelfont={bf,sf},labelsep=period,skip=6pt]{caption}
\usepackage{url}

\definecolor{accent}{HTML}{1B4F72}
\definecolor{warm}{HTML}{C0562A}
\definecolor{slate}{HTML}{56636E}
\definecolor{ink}{HTML}{1E1E1E}

\usepackage[colorlinks=true,linkcolor=accent,citecolor=accent,urlcolor=accent]{hyperref}

% Lean identifiers: typewriter, breakable at dots and underscores.
\DeclareUrlCommand\lean{\urlstyle{tt}}

% ---------- layout ----------
\linespread{1.06}
\setlength{\parskip}{0.55em plus 0.1em minus 0.05em}
\setlist{itemsep=0.25em, topsep=0.2em, parsep=0.1em}
\titleformat{\section}{\sffamily\bfseries\Large\color{accent}}{\thesection}{0.7em}{}
\titleformat{\subsection}{\sffamily\bfseries\normalsize\color{ink}}{\thesubsection}{0.6em}{}
\titleformat{\paragraph}[runin]{\sffamily\bfseries\small\color{accent}}{}{0pt}{}[\enspace]
\titlespacing*{\section}{0pt}{22pt}{5pt}
\titlespacing*{\subsection}{0pt}{13pt}{1pt}
\titlespacing*{\paragraph}{0pt}{9pt}{0pt}
\pagestyle{fancy}
\fancyhf{}
\renewcommand{\headrulewidth}{0pt}
\fancyhead[L]{\sffamily\footnotesize\color{slate}\paperrunning}
\fancyhead[R]{\sffamily\footnotesize\color{slate}\thepage}
\setlength{\headheight}{14pt}
\setcounter{tocdepth}{1}
\setlength{\cftbeforesecskip}{2pt}
\renewcommand{\cftsecfont}{\small}
\renewcommand{\cftsecpagefont}{\small\color{slate}}
\renewcommand{\cftsecleader}{\cftdotfill{\cftdotsep}}
\setlength{\cftsecnumwidth}{1.7em}
\setlength{\emergencystretch}{2em}
\numberwithin{equation}{section}

% ---------- title block ----------
\newcommand{\papertitle}{}
\newcommand{\papersubtitle}{}
\newcommand{\paperrunning}{}
\newcommand{\paperauthors}{KNT\enspace\textperiodcentered\enspace Astra\enspace\textperiodcentered\enspace Fable~5.2}
\newcommand{\paperdate}{Draft of 11 November 2026}
\newcommand{\maketitleblock}{%
  \thispagestyle{plain}%
  {\sffamily\bfseries\color{accent}\fontsize{25}{29}\selectfont\raggedright\papertitle\par}%
  \vspace{5pt}%
  {\sffamily\Large\color{slate}\raggedright\papersubtitle\par}%
  \vspace{8pt}%
  {\color{accent}\hrule height 1.2pt}%
  \vspace{7pt}%
  {\sffamily\small\color{ink}\paperauthors\hfill\color{slate}\paperdate\par}%
  \vspace{10pt}}

% Small headings inside the front matter.
\newcommand{\fronthead}[1]{\par\vspace{6pt}{\sffamily\bfseries\normalsize\color{ink}#1}\par\vspace{-2pt}}
\makeatletter
\newcommand{\contentsblock}{%
  \par\noindent\begin{minipage}{\linewidth}%
  \fronthead{Contents}%
  \setlength{\parskip}{0pt}\setlength{\columnsep}{2.2em}\raggedright
  \begin{multicols}{2}
  \@starttoc{toc}
  \end{multicols}
  \end{minipage}\par}
\makeatother

% ---------- statements ----------
% Theorems, lemmas, propositions and corollaries share one counter.
% Each statement is in a tinted box with an accent rule at the left.
\newtheoremstyle{boxed}{0pt}{0pt}{\normalfont}{0pt}{\sffamily\bfseries\small\color{accent}}{.}{0.6em}{}
\theoremstyle{boxed}
\newtheorem{theorem}{Theorem}[section]
\newtheorem{lemma}[theorem]{Lemma}
\newtheorem{proposition}[theorem]{Proposition}
\newtheorem{corollary}[theorem]{Corollary}

\tcbset{guidebox/.style={enhanced, parbox=false, boxrule=0pt, frame hidden,
  arc=1.5pt, left=11pt, right=10pt, top=7pt, bottom=7pt, before skip=11pt, after skip=9pt}}
% Not breakable: a label at the start of a box that breaks at the top of a page can be lost.
\tcbset{statementbox/.style={guidebox, colback=accent!6,
  borderline west={3pt}{0pt}{accent}}}
\tcolorboxenvironment{theorem}{statementbox}
\tcolorboxenvironment{lemma}{statementbox}
\tcolorboxenvironment{proposition}{statementbox}
\tcolorboxenvironment{corollary}{statementbox}

% The main result of a paper: a framed box with a title bar.
% \begin{maintheorem}{name} uses the counter of the theorems, and \label works after it.
\newtcolorbox{mainresult}[1]{enhanced, parbox=false, colback=white, colframe=accent, boxrule=0.9pt,
  arc=1.5pt, left=11pt, right=11pt, top=8pt, bottom=8pt, title={#1},
  fonttitle=\sffamily\bfseries\small, colbacktitle=accent, coltitle=white,
  before skip=12pt, after skip=12pt}
\newenvironment{maintheorem}[1]%
  {\refstepcounter{theorem}\begin{mainresult}{Theorem~\thetheorem\ (#1)}}%
  {\end{mainresult}}

% ---------- proofs ----------
% A proof has a thin rule at the left, so that the eye finds the statements.
\renewcommand{\proofname}{\upshape\sffamily\bfseries\small\color{slate}Proof}
\renewcommand{\qedsymbol}{\textcolor{slate}{\rule{0.55em}{0.55em}}}
\tcolorboxenvironment{proof}{enhanced, breakable, parbox=false, boxrule=0pt, frame hidden,
  colback=white, borderline west={1.2pt}{0pt}{slate!35}, left=11pt, right=0pt, top=1pt, bottom=2pt,
  before skip=6pt, after skip=11pt}
% A step or a case inside a proof.
\newcommand{\step}[1]{{\sffamily\bfseries\small\color{slate}#1}}

% ---------- idea, notation and fine-print boxes ----------
\newtcolorbox{idea}[1]{guidebox, breakable, colback=warm!6, borderline west={3pt}{0pt}{warm!80},
  before upper={{\sffamily\bfseries\footnotesize\color{warm!75!black}IDEA\enspace\textbar\enspace}{\sffamily\bfseries\small #1}\par}}
\newtcolorbox{notation}[1]{guidebox, breakable, colback=black!4, borderline west={3pt}{0pt}{slate!60},
  fontupper=\small,
  before upper={{\sffamily\bfseries\footnotesize\color{slate}NOTATION\enspace\textbar\enspace}{\sffamily\bfseries\small #1}\par}}
\newtcolorbox{fineprint}{guidebox, breakable, colback=white, borderline west={1.5pt}{0pt}{black!25},
  fontupper=\footnotesize\color{black!78}, left=9pt, top=1pt, bottom=1pt,
  before upper={{\sffamily\bfseries Fine print.}\enspace}}

% ---------- figures ----------
\tikzset{
  vtx/.style={circle, fill=ink, inner sep=1.35pt},
  port/.style={circle, draw=accent, fill=white, line width=0.8pt, inner sep=1.5pt},
  edge/.style={line width=0.55pt, draw=ink},
  hot/.style={line width=1.1pt, draw=warm},
  cool/.style={line width=1.1pt, draw=accent},
  region/.style={draw=accent, fill=accent!7, line width=0.6pt, rounded corners=6pt},
  note/.style={font=\sffamily\footnotesize, text=slate},
}

% ---------- macros of the mathematics ----------
\newcommand{\Z}{\mathcal{Z}}
\newcommand{\PP}{\mathbb{P}}
\newcommand{\EE}{\mathbb{E}}
\newcommand{\comp}{\operatorname{comp}}
\newcommand{\tr}{\operatorname{tr}}
\newcommand{\lgk}[1]{\log^{(#1)}}
\newcommand{\lgs}{\log^{*}}
\newcommand{\sm}{\setminus}
\newcommand{\half}{\tfrac12}
\newcommand{\cO}{\mathcal{O}}


\renewcommand{\papertitle}{The large-girth theorem\\ with one family of short cycles}
\renewcommand{\papersubtitle}{A research note on Section 8 of ``The maximum number of cycles in a graph of large rank''}
\renewcommand{\paperrunning}{The large-girth theorem with one family of short cycles}
\renewcommand{\paperauthors}{Fable~5.2, for KNT}
\renewcommand{\paperdate}{Research note of 12 November 2026}
\hypersetup{pdftitle={The large-girth theorem with one family of short cycles},pdfauthor={Fable 5.2, for KNT}}

% A statement of the paper [M], in the numbering of the draft of 42 pages.
\newcommand{\M}[1]{[M,~#1]}

\begin{document}
\maketitleblock

This note answers the research task of 12 November 2026: find a shorter mathematical proof of
Theorem~8.1 of the paper~[M], the theorem for large girth. The numbers of the statements of~[M]
are those of the draft of 42~pages. The note gives two new proofs. Each proof deletes one maximal
family of short cycles, at one scale. Each proof removes Theorem~8.8 (raising the girth), the
nested sets $S_0\subseteq S_1\subseteq\cdots$, the two girth scales of Section~8.4, and the
deleted set in Propositions~8.5 and~8.7. The manuscript is not changed.

\fronthead{The result in seven sentences}
\begin{enumerate}[leftmargin=1.6em]
\item In the present proof a block of a chain is a piece only if it avoids a deleted set of short
cycles. So the proof needs a small deleted set and a large one, and Theorem~8.8 constructs both.
\item It is sufficient to discard the blocks that \emph{contain} a short cycle. These blocks are
disjoint. So their number is at most a packing number: the largest number of pairwise disjoint
cycles of length less than $C_0\log(\Lambda+2)$.
\item With this change, one proposition on a weighted family of regions gives the packing bound
for short cycles and also the final contradiction for cycles of moderate length
(Proposition~\ref{prop:heavy}).
\item The packing bound for the length $d$ then follows from the packing bound for the length
$C_0\log(d+2)$, by induction. Nothing is deleted, and Theorem~8.8 is not necessary
(Section~\ref{sec:first}).
\item A second proof moves this induction into the lemma on pieces. The new lemma says: at a
fixed girth $D_*$, every large connected subcubic multigraph with a small boundary has a forest
proportion of less than $\frac35$. Then two large disjoint regions of $G-P$ with small boundaries
always give $q(G,P)\le\frac9{25}$ (Section~\ref{sec:second}).
\item Measured in copies of the paper, Section~8 has 271 typeset lines now, 259 lines with the
first proof and 289 lines with the second proof. The paper keeps 42 pages. So the gain is in the
structure of the proof and not in its length.
\item The hypotheses of the final application do not help (Section~\ref{sec:p3}), and each
proof keeps an induction over about $\log^{*}r$ scales (Section~\ref{sec:stays}).
\end{enumerate}

\begin{notation}{the status of this note}
\begin{itemize}[leftmargin=1.2em]
\item \emph{Established on paper.} Sections~\ref{sec:first} and~\ref{sec:second} contain complete
proofs. I checked each step against the statements of~[M]. Two copies of the paper with the new
Section~8 build without an undefined reference.
\item \emph{Tested by computer.} Scripts test the facts on joins of Section~\ref{sec:joins} and
the examples of Sections~\ref{sec:p3} and~\ref{sec:stays} on small graphs
(Section~\ref{sec:checks}). These tests do not replace a proof.
\item \emph{Not done.} There is no formal proof of the new arguments. The Lean proof of~[M]
follows the present Section~8. No other person or model has read this note.
\item \emph{The manuscript.} It is unchanged. Two scripts in the folder of this note apply the
first or the second proof to a copy of the paper.
\end{itemize}
\end{notation}

\contentsblock

\section{The setting}\label{sec:setting}

\subsection{Notation and the theorem}

Throughout, $G$ is a connected cubic multigraph, $\xi$ is a uniformly random even subgraph of
$G$, and $q(G,P)$ is the probability that $\xi\cap E(G-P)$ is a forest. A \emph{region} is a
nonempty set $A$ of vertices such that $G[A]$ is connected. It is \emph{cyclic} if $G[A]$ contains
a cycle. The set $\delta A$ consists of the edges with exactly one end in $A$; these are the
\emph{boundary edges} of $A$. For disjoint regions $A$ and $B$, a \emph{link} is an edge between
them, or a component of $G-(A\cup B)$ that is joined to both. Their number is $\ell(A,B)$. The
event $\cO_A$ is $\xi\cap\delta A=\emptyset$. Logarithms have base $2$, and $\ln$ is the natural
logarithm.

\begin{theorem}[{\M{Theorem~8.1}}]\label{thm:half}
There are integers $D_*\ge7$ and $r_5$ with the following property. Let $G$ be a connected cubic
multigraph of rank $r\ge r_5$. Let $P$ be a nonempty set of at most $3\log r$ vertices such that
$G[P]$ is connected. If $G-P$ has girth at least $D_*$, then $q(G,P)<\frac25$.
\end{theorem}

\subsection{The present proof, and the place of the two girth scales}\label{sec:present}

The proof of~[M] supposes that $q(G,P)\ge\frac25$. It has five steps.
\begin{enumerate}[leftmargin=1.6em]
\item \emph{Pieces} (Lemmas~8.2 and~8.3). Let $Y'\subseteq V(G)\sm P$ be a connected set such
that $G[Y']$ has at most $2\Lambda$ boundary vertices, at least $C_1\Lambda^{4}$ vertices, and
girth at least $C_0\log(\Lambda+2)$. Then $\xi\cap E(G[Y'])$ is a forest with a probability of
less than $\frac35$, for every choice of the edges outside it. Two disjoint sets of this kind give
$q(G,P)\le\frac9{25}$.
\item \emph{A chain} (Lemma~8.4, Proposition~8.5). Pairwise disjoint regions that separate the
same two vertices are nested, and the blocks between them have small boundaries. A block must have
a large girth. For this purpose the proposition takes a set $U\supseteq P$ such that $G-U$ has
large girth, and it discards the blocks that meet $U\sm P$. So $|U\sm P|$ must be smaller than the
number of blocks.
\item \emph{Packing} (Lemma~8.6, Proposition~8.7). If $G-(P\cup S)$ has girth at least $D$, then
it contains at most $C_2\,d^{6}(|S|+1)\,2^{d}$ pairwise disjoint cyclic regions with at most
$d\le2D$ boundary edges.
\item \emph{Raising the girth} (Theorem~8.8). Step~3 for $D=D_*,2D_*,4D_*,\dots$ gives nested
sets $S_0\subseteq S_1\subseteq\cdots$ such that $G-(P\cup S_j)$ has girth greater than
$D_j=2^{j}D_*$ and $\log(|S_j|+1)\le2D_j+C_3\log^{2}D_j$.
\item \emph{Two scales} (Section~8.4). Put $x=\log r$. The second moment of the cycles of
moderate length needs the scale $D^{+}\approx x/16$, and there $|S^{+}|=r^{1/4+o(1)}$. The chain
of Step~2 has about $r^{1/32}$ members. So its blocks cannot avoid $S^{+}$. The proof takes the
girth of the blocks from a second scale $D^{-}\approx\sqrt x$, where $|S^{-}|=r^{o(1)}$.
\end{enumerate}
So the two scales have one cause: Step~2 obtains the girth of a block from a deleted set, and the
deleted set of the large scale is larger than the chain. Section~\ref{sec:first} removes this
cause.

\subsection{Statements of [M] that the note uses}\label{sec:used}

\begin{fineprint}
\textbf{(a) \M{Lemma~2.4} and its proof.} Let $X_1,\dots,X_n$ be pairwise disjoint cyclic regions
of a connected multigraph with $|\delta X_i|\le d$. Let $E_i$ be the event that $\cO_{X_i}$ holds
and $\xi\cap E(G[X_i])\ne\emptyset$, and put $p_i=\PP(E_i)$. Then $p_i\ge2^{-d-1}$, and
$\PP(E_i\cap E_j)=2^{\ell(X_i,X_j)-1}p_ip_j$ for $i\ne j$. For each $i$ at most $d$ indices $j$
have $\ell(X_i,X_j)\ge3$. The probability that no $E_i$ occurs is at most
$\half+(d+1)2^{d-1}/n$.

\smallskip
\textbf{(b) \M{Section~3.1 and Proposition~3.5}.} Let $U$ be a set of vertices such that $J=G-U$
is a simple graph of girth greater than $D$. Put $s=\lfloor(D-1)/2\rfloor$. Let $\mathcal C$ be
the set of the cycles of $J$ of length at most $L$, with the weights $a_C=2^{-|C|}$ of total
$\mu$. Then, for every vertex $w$,
\begin{equation}\label{eq:load}
  \sum_{C\in\mathcal C,\ C\ni w}a_C\;\le\;L\,2^{-s}.
\end{equation}
There is a random variable $Z\ge0$ with $\EE Z=\mu$ such that $\xi$ contains a cycle of $J$ if
$Z\ne0$. Let $W_2$ be the sum of $a_Ca_{C'}$ over the ordered pairs of disjoint cycles of
$\mathcal C$ whose vertex sets have exactly two links. Proposition~3.5 says the following. Let $r$
be large, $x=\log r$ and $y=\log x$. Let $P$ be a set of at most $2r/x^{2}$ vertices such that
$G[P]$ has at most $x/100$ components, and let $q(G,P)\ge\frac25$. Let $D$ and $L$ be integers
with $x/16\le D\le x/8$ and $7x\le L\le x\sqrt y$. Let $S$ be the union of the vertex sets of a
maximal family of pairwise disjoint cycles of length at most $D$ in $G-P$, and let
$|S|\le r/x^{2}$. Then, for $U=P\cup S$,
\[
  L^{2}2^{-s}\le r^{-1/60},\qquad \mu\ge\frac14\ln\frac{L}{6x}-\frac1{\sqrt x},\qquad
  \EE Z^{2}\le\mu^{2}+W_2+\bigl(1+r^{-1/60}\bigr)\mu .
\]

\smallskip
\textbf{(c) \M{Lemma~8.2}.} There are absolute constants $C_0$ and $C_1\ge1$ with the following
property. Let $\Lambda\ge4$, and let $N$ be a connected subcubic multigraph with at most
$2\Lambda$ boundary vertices, at least $C_1\Lambda^{4}$ vertices, and girth at least
$C_0\log(\Lambda+2)$. Then every set of boundary vertices of even size has forest proportion less
than $\frac35$. (The boundary vertices are the vertices of degree less than three. For an even
set $T$ of them, a $T$-join is a set of edges in which exactly the vertices of $T$ have odd
degree, and the forest proportion of $T$ is the proportion of the $T$-joins that are forests.)

\smallskip
\textbf{(d) \M{Lemma~8.3}.} Let $U_1,\dots,U_m\subseteq V(G)\sm P$ be pairwise disjoint sets such
that each $G[U_i]$ is connected and every even set of boundary vertices of $G[U_i]$ has forest
proportion at most $\frac35$. Then $q(G,P)\le(3/5)^{m}$.

\smallskip
\textbf{(e) \M{Lemma~8.4}.} Let $u$ and $p$ be two vertices of a connected graph $H$. Pairwise
disjoint connected sets that avoid $u$ and $p$ and separate $u$ from $p$ are linearly ordered:
$C_1<C_2$ means that $C_1\cup U_1\subseteq U_2$, where $U_i$ is the vertex set of the component of
$u$ in $H-C_i$.

\smallskip
\textbf{(f) \M{Lemma~8.6}.} Let $p$ be a vertex of a connected multigraph. Let $\mathcal A$ be a
finite family of regions that avoid $p$, with weights $w_A\ge0$ of total $\mu>0$. Let $W_2$ be the
sum of $w_Aw_{A'}$ over the ordered pairs of disjoint regions of $\mathcal A$ with exactly two
links. Let $N$ and $e$ be numbers such that, for every $A\in\mathcal A$, at most $N$ vertices
outside $A$ have a neighbour in $A$, and the regions of $\mathcal A$ that are disjoint from $A$
and joined to $A$ by an edge have total weight at most $e$. Let $K$ be a number with
$W_2\ge\frac23\mu^{2}-K\mu$ and $\mu>\frac32K+3e$. Then there are a region $A\in\mathcal A$ and a
vertex $u\notin A$ in the component of $p$ in $G-A$ such that the regions of $\mathcal A$ that
are disjoint from $A$, avoid $u$, and separate $u$ from $p$ in $G-A$ have total weight at least
$(\mu/3-K/2-e)/N$.
\end{fineprint}

The next lemma is the step ``Blocks'' in the proof of \M{Proposition~8.5}. Both new proofs use
it, so the note states it alone.

\begin{lemma}[Blocks]\label{lem:blocks}
Let $u$ and $p$ be two vertices of a connected graph $H$. Let $C_1<C_2<\dots<C_n$ be pairwise
disjoint connected sets that avoid $u$ and $p$ and separate $u$ from $p$, and let $w\ge1$ be an
integer. Then there are pairwise disjoint sets $Y_l\subseteq V(H)$, one for each integer $l\ge0$
with $(l+1)w\le n-1$, such that
\begin{enumerate}
\item[\textup{(i)}] $Y_l$ contains $w-1$ of the sets $C_i$;
\item[\textup{(ii)}] every edge of $H$ that leaves $Y_l$ ends in $C_{lw+1}$ or in $C_{(l+1)w+1}$;
\item[\textup{(iii)}] $Y_l$ lies in the component of $u$ in $H-C_{(l+1)w+1}$.
\end{enumerate}
\end{lemma}

\begin{proof}
Let $U_i$ be the vertex set of the component of $u$ in $H-C_i$. Then $U_i\cup C_i\subseteq U_j$
for $i<j$. For $i<j$ put $X_{ij}=U_j\sm(U_i\cup C_i)$. An edge of $H$ that leaves $U_j$ ends
in $C_j$, because $U_j$ is a component of $H-C_j$. No edge joins $X_{ij}$ to $U_i$: its end in
$X_{ij}$ would lie in the component $U_i$ of $H-C_i$. So every edge of $H$ that leaves $X_{ij}$
ends in $C_i$ or in $C_j$. If $i<l<j$, then $C_l\subseteq U_j$, and $C_l$ avoids
$U_i\cup C_i\subseteq U_l$; so $C_l\subseteq X_{ij}$. If $i<j\le i'<j'$, then
$X_{ij}\subseteq U_j\subseteq U_{i'}$ is disjoint from $X_{i'j'}$. So the sets
$Y_l=X_{lw+1,\,(l+1)w+1}$ have the three properties.
\end{proof}

\section{The first proof: count the blocks that contain a short cycle}\label{sec:first}

Let $C_0$ and $C_1$ be the constants of \M{Lemma~8.2}. We may suppose that $C_0\ge1$, because a
larger number $C_0$ makes the hypothesis of that lemma stronger.

Call a cycle \emph{$\Lambda$-short} if its length is less than $C_0\log(\Lambda+2)$. Let
$\nu_\Lambda(G,P)$ be the largest number of pairwise disjoint $\Lambda$-short cycles of $G-P$. Let
$D_*\ge7$ be a fixed integer such that every integer $d\ge D_*$ satisfies
\begin{equation}\label{eq:dstar}
  \bigl(C_0\log(d+2)\bigr)^{C_0+9}\le d,\qquad (d+2)^{C_0}\le2\,d^{\,C_0},\qquad d^{2}\ge20\,C_1 .
\end{equation}
Then $C_0\log(d+2)\le d$. This number $D_*$ replaces the number $D_*$ of~[M], which satisfies
$D\ge C_0\log(2D+2)$ for $D\ge D_*$.

\begin{proposition}[A long chain]\label{prop:chain}
Let $G$ be a connected cubic multigraph, and let $P$ be a nonempty set of vertices such that
$G[P]$ is connected. Let $p\in P$, let $A$ be a region that avoids $P$, and let $u\notin A$ be a
vertex in the component of $p$ in $G-A$. Let $C_1,\dots,C_T$ be pairwise disjoint regions that
avoid $P$, $A$ and $u$ and separate $u$ from $p$ in $G-A$. Let $\Lambda\ge4$ be an integer and
$m\ge1$ a number such that $A$ and every $C_i$ have at most $\Lambda$ boundary edges, and every
$C_i$ has at least $m$ vertices. If
\[
  T\;>\;\bigl(\Lambda+\nu_\Lambda(G,P)+2\bigr)\bigl(\lceil2C_1\Lambda^{5}/m\rceil+1\bigr),
\]
then $q(G,P)\le\frac9{25}$.
\end{proposition}

\begin{proof}
Each $C_i$ is connected and separates $u$ from $p$ in $G-A$. So it lies in the component $H$ of
$G-A$ that contains $u$ and $p$. By \M{Lemma~8.4} we may assume that $C_1<C_2<\dots<C_T$. Put
$w=\lceil2C_1\Lambda^{5}/m\rceil+1$ and $\nu=\nu_\Lambda(G,P)$.

\step{Blocks.} Lemma~\ref{lem:blocks} gives at least $\Lambda+\nu+2$ pairwise disjoint sets
$Y_l$. Each of them contains $w-1$ of the sets $C_i$, so $|Y_l|\ge(w-1)m\ge2C_1\Lambda^{5}$. Every
edge of $H$ that leaves $Y_l$ ends in one of two sets $C_i$, so at most $2\Lambda$ edges of $H$
leave $Y_l$. The set $P$ is connected, it contains $p$, and it is disjoint from $A$ and from the
sets $C_i$. So $P$ lies in the component of $p$ in $H-C_i$ for every $i$, and every $Y_l$ avoids
$P$, by~(iii).

Discard $Y_l$ if it contains a vertex with a neighbour in $A$, or if $G[Y_l]$ contains a
$\Lambda$-short cycle. The set $A$ has at most $\Lambda$ boundary edges, and the sets $Y_l$ are
pairwise disjoint subsets of $V(G)\sm P$. So this discards at most $\Lambda+\nu$ sets, and two
sets remain. Let $Y$ be one of them. The edges of $G$ that leave $Y$ are edges of $H$, so
$|\delta Y|\le2\Lambda$, and $G[Y]$ has girth at least $C_0\log(\Lambda+2)$.

\step{Pieces.} Every component of $G[Y]$ has an edge that leaves $Y$. So there are at most
$2\Lambda$ components, and one of them has a vertex set $Y'$ with at least $C_1\Lambda^{4}$
vertices. A boundary vertex of $G[Y']$ is an end of an edge of $\delta Y$, so there are at most
$2\Lambda$ of them. By \M{Lemma~8.2}, every even set of boundary vertices of $G[Y']$ has forest
proportion less than $\frac35$. \M{Lemma~8.3}, for the two sets $Y'$, gives
$q(G,P)\le\frac9{25}$.
\end{proof}

The difference from \M{Proposition~8.5} is in the second paragraph of the step ``Blocks''. There
is no set $U$ and no number $D$, and the conclusion is for the marked set $P$.

\begin{proposition}[Heavy families]\label{prop:heavy}
Let $G$ be a connected cubic multigraph, and let $P$ be a nonempty set of vertices such that
$G[P]$ is connected and $q(G,P)\ge\frac25$. Let $\mathcal A$ be a finite family of regions that
avoid $P$, with weights $w_A\ge0$ of total $\mu$, and let $W_2$ be the sum of $w_Aw_{A'}$ over the
ordered pairs of disjoint regions of $\mathcal A$ with exactly two links. Let $\Lambda\ge4$ be an
integer, and let $m\ge1$, $\kappa>0$ and $K\ge0$ be numbers such that
\begin{enumerate}
\item[\textup{(a)}] every region of $\mathcal A$ has at most $\Lambda$ boundary edges and at
least $m$ vertices;
\item[\textup{(b)}] for every $A\in\mathcal A$, the regions of $\mathcal A$ that meet $A$ have
total weight at most $\kappa$;
\item[\textup{(c)}] some random variable $Z\ge0$ has $\EE Z=\mu$ and
$\EE Z^{2}\le\mu^{2}+W_2+K\mu$, and $Z=0$ whenever $\xi\cap E(G-P)$ is a forest.
\end{enumerate}
Then
\begin{equation}\label{eq:heavy}
  \mu\;\le\;\tfrac32K+3\Lambda\kappa\,\bigl(\Lambda+\nu_\Lambda(G,P)+3\bigr)
  \bigl(\lceil2C_1\Lambda^{5}/m\rceil+1\bigr).
\end{equation}
\end{proposition}

\begin{proof}
Put $\nu=\nu_\Lambda(G,P)$ and $w=\lceil2C_1\Lambda^{5}/m\rceil+1$. We may assume that
$\mu>\frac32K+3\Lambda\kappa$. By~(c), $\PP(Z\ne0)\le1-q(G,P)\le\frac35$, and the Cauchy--Schwarz
inequality $\mu^{2}\le\EE Z^{2}\cdot\PP(Z\ne0)$ gives $W_2\ge\frac23\mu^{2}-K\mu$.

Fix $p\in P$. We apply \M{Lemma~8.6}. Let $A\in\mathcal A$. By~(a), at most $\Lambda$ vertices
outside $A$ have a neighbour in $A$, and a region that is joined to $A$ by an edge contains one of
them. The regions of $\mathcal A$ that contain a given vertex all meet one of them, so their total
weight is at most $\kappa$ by~(b). Hence we may take $N=\Lambda$ and $e=\Lambda\kappa$. The lemma
gives a region $A\in\mathcal A$, a vertex $u\notin A$ in the component of $p$ in $G-A$, and
regions of $\mathcal A$ of total weight at least $(\mu/3-K/2-\Lambda\kappa)/\Lambda$ that are
disjoint from $A$, avoid $u$, and separate $u$ from $p$ in $G-A$. Take a maximal family
$C_1,\dots,C_T$ of pairwise disjoint regions among them. Each of these regions meets some $C_i$,
so their total weight is at most $T\kappa$ by~(b). As $q(G,P)>\frac9{25}$,
Proposition~\ref{prop:chain} gives $T\le(\Lambda+\nu+2)\,w$. So
$\mu/3-K/2-\Lambda\kappa\le\Lambda\kappa\,(\Lambda+\nu+2)\,w$, and the claim follows because
$w\ge1$.
\end{proof}

In other words: if a family with the properties (a), (b) and (c) has a total weight above the
bound~\eqref{eq:heavy}, then $q(G,P)<\frac25$. The proof has two cases. Either the pairs with two
links have a small weight, and the second moment gives $q(G,P)<\frac25$. Or they give two large
disjoint regions with small boundaries and without a short cycle, and the product gives
$q(G,P)\le\frac9{25}$.

\begin{proposition}[Packing]\label{prop:packing}
Let $G$ be a connected cubic multigraph, and let $P$ be a nonempty set of vertices such that
$G[P]$ is connected. Suppose that $q(G,P)\ge\frac25$ and that $G-P$ has girth at least $D_*$. Let
$d\ge D_*$ be an integer. Then
\begin{enumerate}
\item[\textup{(i)}] $G-P$ has at most $2^{d}d^{\,C_0+9}$ pairwise disjoint cycles of length at
most $d$;
\item[\textup{(ii)}] $\nu_d(G,P)\le2\,d^{\,C_0+1}$.
\end{enumerate}
\end{proposition}

\begin{proof}
We use induction on $d$. Suppose that~(i) holds for every smaller integer that is at least $D_*$.

\step{The shorter cycles.} Let $g$ be the largest integer that is less than $C_0\log(d+2)$. Then
$g<d$ by~\eqref{eq:dstar}, and a cycle is $d$-short if and only if its length is at most $g$. If
$g<D_*$, then $G-P$ has no such cycle. Otherwise~(i) for $g$ and~\eqref{eq:dstar} give
\[
  \nu_d(G,P)\;\le\;2^{g}g^{\,C_0+9}\;\le\;(d+2)^{C_0}\bigl(C_0\log(d+2)\bigr)^{C_0+9}
  \;\le\;2\,d^{\,C_0+1}.
\]
This is~(ii).

\step{The cycles of length at most $d$.} Let $Q_1,\dots,Q_n$ be pairwise disjoint cycles of
length at most $d$ in $G-P$, and put $X_i=V(Q_i)$. Every vertex of $Q_i$ lies on exactly one edge
outside $Q_i$, and $Q_i$ has at least $D_*$ vertices. So $X_i$ is a cyclic region with at most
$d$ boundary edges and at least $D_*$ vertices. Let $E_i$ be the event that $\cO_{X_i}$ holds and
$\xi\cap E(G[X_i])\ne\emptyset$. Put $p_i=\PP(E_i)$, let $Z$ be the number of the events $E_i$
that occur, and put $\mu=\EE Z$. By Section~\ref{sec:used}(a), $p_i\ge2^{-d-1}$ and
$\PP(E_i\cap E_j)=2^{\ell(X_i,X_j)-1}p_ip_j$ for $i\ne j$, and for each $i$ at most $d$ indices
$j$ have $\ell(X_i,X_j)\ge3$. Let $W_2$ be the sum of $p_ip_j$ over the ordered pairs $(i,j)$ with
$\ell(X_i,X_j)=2$. Then $\EE Z^{2}\le\mu^{2}+W_2+(d+1)\mu$: a pair with at most one link gives at
most $p_ip_j$, a pair with two links gives $2p_ip_j$, a pair with three or more links gives at
most $p_i$, and the terms with $i=j$ give $\mu$. If $\xi\cap E(G-P)$ is a forest, then no event
$E_i$ occurs, and $Z=0$.

So Proposition~\ref{prop:heavy} applies to the sets $X_i$ with the weights $p_i$, and with
$\Lambda=d$, $m=D_*$, $\kappa=1$ and $K=d+1$; property~(b) holds because the sets are pairwise
disjoint and $p_i\le1$. We have $\lceil2C_1d^{5}/D_*\rceil+1\le C_1d^{5}$, because $D_*\ge7$, and
$d+3\le d^{\,C_0+1}$. With~(ii),
\[
  \mu\;\le\;\tfrac32(d+1)+3d\,\bigl(d+3+2\,d^{\,C_0+1}\bigr)\,C_1d^{5}\;\le\;10\,C_1\,d^{\,C_0+7}.
\]
Hence $n\le2^{d+1}\mu\le20\,C_1\,2^{d}d^{\,C_0+7}\le2^{d}d^{\,C_0+9}$, by~\eqref{eq:dstar}. This
is~(i).
\end{proof}

\subsection{Proof of Theorem~\ref{thm:half}}\label{sec:firstproof}

Let $G$ and $P$ be as in the theorem, with the number $D_*$ of~\eqref{eq:dstar}, and suppose that
$q(G,P)\ge\frac25$. We derive a contradiction when $r$ is large. All limits are for
$r\to\infty$, and they are uniform in $G$ and $P$. Put
\[
  x=\log r,\qquad A=6\,e^{124},\qquad L=\lfloor Ax\rfloor,\qquad D=\lceil x/16\rceil .
\]

\step{Cycles of moderate length.} Let $S$ be the union of the vertex sets of a maximal family of
pairwise disjoint cycles of length at most $D$ in $G-P$. By
Proposition~\ref{prop:packing}(i), $|S|\le2^{D}D^{\,C_0+10}\le r/x^{2}$ for large $r$. We use
\M{Proposition~3.5} for $G$, $P$, $D$, $L$ and $S$. Its hypotheses hold for large $r$: the set
$P$ has at most $3x$ vertices and $G[P]$ is connected, $x/16\le D\le x/8$, and
$7x\le L\le x\sqrt{\log x}$. So we are in the setting of Section~\ref{sec:used}(b), with
$U=P\cup S$ and $J=G-U$. With $\eta=r^{-1/60}$ the proposition gives
\[
  L^{2}2^{-s}\le\eta,\qquad \mu\;\ge\;\frac14\ln\frac{L}{6x}-\frac1{\sqrt x}\;>\;30,\qquad
  \EE Z^{2}\;\le\;\mu^{2}+W_2+(1+\eta)\mu .
\]
The bound $\mu>30$ holds because $L/(6x)\ge e^{124}-1$.

\step{The contradiction.} We apply Proposition~\ref{prop:heavy} to the vertex sets of the cycles
of $\mathcal C$, with the weights $a_C$, and with $\Lambda=L$, $m=x/16$, $\kappa=\eta$ and
$K=1+\eta$. A cycle $C\in\mathcal C$ has at most $L$ vertices, each of them has one edge end
outside $C$, and the girth of $J$ is greater than $D\ge x/16$. So $V(C)$ has at most $L$ boundary
edges and at least $x/16$ vertices; this is~(a). A cycle of $\mathcal C$ that meets $C$ passes
through one of the at most $L$ vertices of $C$, so~(b) holds by~\eqref{eq:load}. If
$\xi\cap E(G-P)$ is a forest, then $\xi$ contains no cycle of $J$, and $Z=0$; this is~(c). By
Proposition~\ref{prop:packing}(ii), $\nu_L(G,P)\le2L^{C_0+1}$. So
\[
  \mu\;\le\;\tfrac32(1+\eta)+3L\eta\,\bigl(L+2L^{C_0+1}+3\bigr)
  \bigl(\lceil32C_1L^{5}/x\rceil+1\bigr)\;=\;\tfrac32+o(1),
\]
because $L\le Ax$ and $\eta=r^{-1/60}$. This contradicts $\mu>30$, and
Theorem~\ref{thm:half} is proved. \qed

\section{What the first proof replaces}\label{sec:replaces}

\begin{center}
\small
\begin{tabular}{@{}p{7.1cm}p{7.1cm}@{}}
\toprule
In [M] & In the first proof\\
\midrule
Lemmas 8.2, 8.3, 8.4 and 8.6 & No change.\\\addlinespace
Proposition 8.5 (a long chain): a set $U$ such that $G-U$ has girth at least $D$; it discards the
blocks that meet $U\sm P$; the conclusion is for $q(G,U)$ &
Proposition~\ref{prop:chain}: no set $U$ and no number $D$; it discards the blocks that contain
a short cycle; the conclusion is for $q(G,P)$.\\\addlinespace
The step ``A chain'' in the proof of Proposition 8.7, and the paragraphs ``Pairs with two
links'', ``The chain'' and ``The contradiction'' of Section 8.4. This is one argument, written two
times &
Proposition~\ref{prop:heavy} (heavy families), used two times.\\\addlinespace
Proposition 8.7 (packing): a deleted set $S$, girth at least $D$ outside $P\cup S$, regions with
at most $d\le2D$ boundary edges, bound $C_2\,d^{6}(|S|+1)2^{d}$ &
Proposition~\ref{prop:packing}: no deleted set, every length $d\ge D_*$, bound
$2^{d}d^{\,C_0+9}$, proof by induction on $d$.\\\addlinespace
Theorem 8.8 (raising the girth): the sets $S_0\subseteq S_1\subseteq\cdots$ at the scales
$2^{j}D_*$, with $\log(|S_j|+1)\le2D_j+C_3\log^{2}D_j$ & Removed.\\\addlinespace
Section 8.4: the scales $D^{-}$ and $D^{+}$, the sets $S^{-}\subseteq S^{+}$, the
estimate~(8.1) & One scale $D=\lceil x/16\rceil$ and one set $S$.\\\addlinespace
The number $D_*$ with $D\ge C_0\log(2D+2)$ for $D\ge D_*$ & The number $D_*$
of~\eqref{eq:dstar}. It is larger.\\
\bottomrule
\end{tabular}
\end{center}

\paragraph{The size.} I applied the first proof to a copy of the paper and built the copy. The
count is in lines of the text of the PDF file.
\begin{center}
\small
\begin{tabular}{@{}lccc@{}}
\toprule
 & Now & First proof & Second proof (Section~\ref{sec:second})\\
\midrule
Section 8, typeset lines & 271 & 259 & 289\\
The paper, typeset lines & 1624 & 1615 & 1644\\
The paper, pages & 42 & 42 & 42\\
Numbered statements in Section 8 & 8 & 8 & 9\\
\bottomrule
\end{tabular}
\end{center}
So the first proof is shorter by 12 lines in Section~8. The paper as a whole is shorter by 9
lines, because Section~1 must say that the formal proof raises the girth in stages. This is not a
large change of the length. The change is in the structure: the deleted sets, the two scales and
Theorem~8.8 are no longer in the proof, and the chain argument is written one time.

\paragraph{Remarks.}
\begin{itemize}[leftmargin=1.2em]
\item The new packing bound is better. Theorem~8.8 gives about
$2^{2D}D^{\,C_3\log D}$ vertices in the deleted set of the scale $D$.
Proposition~\ref{prop:packing} gives $2^{D}D^{\,C_0+10}$.
\item The number $D_*$ is larger than before, because~\eqref{eq:dstar} contains the constant
$C_1$. It is an absolute constant, and \M{Theorem~7.1} holds for every fixed bound of the girth.
So the proof of the main theorem of~[M] does not change.
\item The formal proof follows the present Section~8. If the manuscript takes a new proof,
Section~1.4 of~[M] must say that the formal proof of Theorem~8.1 uses another method. The table
in the two copies has this sentence.
\item Proposition~\ref{prop:chain} is used one time, in Proposition~\ref{prop:heavy}, and
\M{Lemma~8.6} is used one time. I did not merge them, because this does not make the proof
shorter.
\end{itemize}

\section{The second proof: pieces at a fixed girth}\label{sec:second}

The first proof has an induction, because \M{Lemma~8.2} needs a girth of about
$C_0\log\Lambda$ for $2\Lambda$ boundary vertices. This section proves the lemma at a fixed
girth. The induction is then in the proof of this lemma, and the other statements have no packing
number.

\subsection{Joins}\label{sec:joins}

Let $N$ be a connected subcubic multigraph. Its \emph{boundary vertices} are the vertices of
degree less than three. Let $T$ be an even set of boundary vertices. A \emph{$T$-join} is a set of
edges in which exactly the vertices of $T$ have odd degree. There are $2^{\beta(N)}$ of them, and
the \emph{forest proportion} of $T$ is the proportion of the $T$-joins that are forests. A
region $A$ of $N$ is \emph{$T$-free} if it contains no vertex of $T$ and every component of $N-A$
contains an even number of vertices of $T$.

\begin{lemma}[Free regions]\label{lem:free}
Let $J$ be a uniformly random $T$-join of $N$. For a region $A$ let $E_A$ be the event that
$J\cap\delta A=\emptyset$ and $J\cap E(N[A])\ne\emptyset$.
\begin{enumerate}
\item[\textup{(i)}] Let $A$ be a region without a vertex of $T$. Some $T$-join avoids $\delta A$
if and only if $A$ is $T$-free.
\item[\textup{(ii)}] If $A$ is $T$-free, then $E_A$ has the same probability as for a uniformly
random even subgraph of $N$.
\item[\textup{(iii)}] If $A$ and $B$ are disjoint $T$-free regions, then $E_A\cap E_B$ is empty,
or it has the same probability as for a uniformly random even subgraph of $N$.
\end{enumerate}
So the bound of Section~\ref{sec:used}(a) holds for $J$ and pairwise disjoint $T$-free cyclic
regions: if there are $n$ of them, each with at most $d$ boundary edges, then the forest
proportion of $T$ is at most $\half+(d+1)2^{d-1}/n$.
\end{lemma}

\begin{proof}
(i) A $T$-join that avoids $\delta A$ has an even number of vertices of odd degree in each
component of $N-A$. In the other direction, each component $K$ of $N-A$ is connected and has a
$(T\cap K)$-join if $|T\cap K|$ is even. The union of these joins is a $T$-join $F$ without an
edge that has an end in $A$.

(ii) The map $J\mapsto J+F$ is a bijection from the $T$-joins to the even subgraphs of $N$, and
it keeps the edges that have an end in $A$. The event $E_A$ depends on these edges only.

(iii) Suppose that some $T$-join is in $E_A\cap E_B$. Delete its edges in $N[A]$ and in $N[B]$.
These form even subgraphs, so the result is a $T$-join $F'$ without an edge that has an end in
$A\cup B$. Use the map $J\mapsto J+F'$.

For the last statement, the proof of \M{Lemma~2.4} uses the probabilities of the events $E_i$ and
upper bounds for the probabilities of the events $E_i\cap E_j$. By (ii) and (iii) these hold
for~$J$. If $J$ is a forest, then no event $E_i$ occurs.
\end{proof}

\begin{lemma}[Product, for joins]\label{lem:product}
Let $N$ be a multigraph, and let $P$ and $T$ be sets of vertices such that $N$ has a $T$-join. Let
$U_1,\ldots,U_m\subseteq V(N)\sm(P\cup T)$ be pairwise disjoint sets of vertices of degree three
such that each $N[U_i]$ is connected. Suppose that in each $N[U_i]$ every even set of boundary
vertices has forest proportion at most $\frac35$. Then $J\cap E(N-P)$ is a forest for at most the
proportion $(3/5)^m$ of the $T$-joins $J$ of $N$.
\end{lemma}

For a cubic multigraph $G$ and $T=\emptyset$ this is \M{Lemma~8.3}. The proof is the proof of
that lemma.

\begin{proof}
Let $E'$ consist of the edges that lie in none of the graphs $N[U_i]$, and let $J$ be a uniformly
random $T$-join. Fix a set $x\subseteq E'$ that extends to a $T$-join. For each $i$, let $T_i$ be
the set of the vertices of $U_i$ that are incident with an odd number of edges of $x$. The set
$U_i$ avoids $T$, so the parity equations say that $J\cap E(N[U_i])$ is a $T_i$-join. These
equations involve disjoint sets of edges, so the joins are independent and uniform, conditional on
$J\cap E'=x$. A vertex of degree three in $N[U_i]$ has no incident edge in $E'$. So $T_i$ is an
even set of boundary vertices. If $J\cap E(N-P)$ is a forest, its restriction to every $N[U_i]$
is a forest. The conditional probability is therefore at most $(3/5)^m$. Average over $x$.
\end{proof}

\subsection{The lemma on pieces}

Let $B$ be the absolute constant of \M{Section~8.1}: each layer in the conclusion of
\M{Lemma~4.1} with $\delta=\frac1{20}$, together with at most $k$ marked vertices, has fewer than
$B(k+1)^{2}$ vertices. For $k\ge1$ let $s_k$ be the least integer with
$2^{s_k}>(2s_k+1)B(k+1)^{2}$. Then $2^{s_k}\le2(2s_k+1)B(k+1)^{2}$, the numbers $s_k$ do not
decrease, and $s_k=O(\log k)$.

\begin{lemma}[Pieces with few short cycles]\label{lem:few}
Let $k\ge2$, and let $N$ be a connected subcubic multigraph with at least two and at most $k$
boundary vertices. Put $s=s_k$. Let $\mathcal D$ be a maximal family of pairwise disjoint cycles
of length at most $2s$ in $N$, and let $t$ be their number. If $N$ has more than
$3\cdot2^{s}(k+2s\,t)$ vertices, then every set of boundary vertices of even size
has forest proportion less than $\frac35$.
\end{lemma}

The lemma has no hypothesis on the girth. For $t=0$ it is \M{Lemma~8.2}: there
$k=2\Lambda$, the girth is greater than $2s$, and $C_1\Lambda^{4}>6\Lambda\,2^{s}$.

\begin{proof}
Fix a boundary vertex $o$, and call the set of the vertices at distance $j$ from $o$ the
\emph{layer} $j$.

\step{A large layer.} A ball of radius $s$ has at most $3\cdot2^{s}$ vertices. So some vertex $w$
has distance more than $s$ from every boundary vertex and from every vertex of a cycle of
$\mathcal D$. Let $W$ be the set of the vertices within distance $s$ of $w$. They have degree
three. Take a breadth-first tree from $w$. Suppose that some vertex at distance less than $s$
from $w$ has fewer than two children in this tree, or that $w$ has fewer than three. Then
$N[W]$ contains a cycle of length at most $2s$. By the maximality of $\mathcal D$, this cycle
meets a cycle of $\mathcal D$, and this contradicts the choice of $w$. So $|W|\ge2^{s}$. The
distances from $o$ to two vertices of $W$ differ by at most
$2s$. So $W$ meets at most $2s+1$ layers, and some layer has at least
$2^{s}/(2s+1)>B(k+1)^{2}$ vertices.

\step{A marker.} This step is the second step of the proof of \M{Lemma~8.2}. Suppose that some
set $T$ of boundary vertices has forest proportion at least $\frac35$. Take a breadth-first tree
of $N$ from $o$, and let $P$ be the union of its paths from $o$ to the other boundary vertices.
Then $N[P]$ is connected, every vertex outside $P$ has degree three, and $P$ has at most $k$
vertices in each layer. The tree paths inside $P$ contain a $T$-join $F_0$. The map
$J\mapsto J+F_0$ is a bijection from the $T$-joins of $N$ to the even subgraphs of $N$, and
$J+F_0$ has the same edges in $N-P$ as $J$. So if $J$ is a forest, then $J+F_0$ is a forest in
$N-P$. Hence $q(N,P)\ge\frac35>\half+\frac1{20}$. The distance from $o$ is a level function. By
\M{Lemma~4.1}, each layer has fewer than $B(k+1)^{2}$ vertices. This contradicts the first step.
\end{proof}

\begin{theorem}[Pieces at a fixed girth]\label{thm:pieces}
There are an integer $D_*\ge7$ and an absolute constant $C_1\ge2$ with the following property. Let
$k\ge1$, and let $N$ be a connected subcubic multigraph of girth at least $D_*$ with at most $k$
boundary vertices and at least $C_1k^{7}$ vertices. Then every set of boundary vertices of even
size has forest proportion less than $\frac35$.
\end{theorem}

\begin{proof}
Fix $k_0$ such that $4s_k<k$ for all $k\ge k_0$, and put $D_*=2s_{k_0}+1$. We define numbers
$F(k)\ge2$ below, and we prove the theorem with $F(k)$ in the place of $C_1k^{7}$, by induction
on $k$.

Let $N$ and $k$ be as in the theorem, with at least $F(k)$ vertices, and put $s=s_k$. If $N$ has
at most one boundary vertex, the only set is the empty one. Then $N$ contains a cycle, because it
has at least two vertices, and the forest proportion is at most $\half$. So let $N$ have at least
two boundary vertices, and let $\mathcal D$ and $t$ be as in Lemma~\ref{lem:few}.

\step{The base.} Let $k<k_0$. Then $2s<D_*$, so $t=0$. By Lemma~\ref{lem:few} the theorem
holds with $F(k)=3\cdot2^{s}k+1$.

\step{The numbers.} Let $k\ge k_0$, and put $k'=4s$. Then $k'<k$. Put
\begin{gather*}
  w=\lceil k'F(k')/D_*\rceil+1,\qquad n_1=5(2s+1)\,2^{2s},\qquad n_2=k^{2}(k+2)\,w,\\
  F(k)=3\cdot2^{s}\bigl(k+2s\,(k+n_1+n_2+2)\bigr)+1 .
\end{gather*}
If $t<k+n_1+n_2+2$, the claim holds by Lemma~\ref{lem:few}. So let $t\ge k+n_1+n_2+2$, and let
$T$ be an even set of boundary vertices. At most $k$ cycles of $\mathcal D$ contain a boundary
vertex. Each of the other cycles $Q$ has at most $2s$ boundary edges, and $V(Q)$ contains no
vertex of $T$. Call $Q$ \emph{free} if $V(Q)$ is $T$-free. So more than $n_1$ cycles are free, or
more than $n_2$ cycles are not free.

\step{More than $n_1$ free cycles.} By Lemma~\ref{lem:free} with $d=2s$, the forest proportion
of $T$ is less than $\half+(2s+1)\,2^{2s-1}/n_1=\frac35$.

\step{More than $n_2$ cycles that are not free.} Let $Q$ be one of them. Some component of
$N-V(Q)$ contains an odd number of vertices of $T$. The set $T$ is even and avoids $Q$. So $T$ has
a vertex $a$ in this component and a vertex $b$ in another component, and $V(Q)$ separates $a$
from $b$. There are fewer than $k^{2}$ pairs $a,b$. So more than $(k+2)w$ of these cycles
separate the same two vertices $a$ and $b$. By \M{Lemma~8.4} their vertex sets are linearly
ordered, and Lemma~\ref{lem:blocks} for $N$, $a$ and $b$ gives $k+2$ pairwise disjoint sets $Y_l$.
At most $k$ of them contain a boundary vertex of $N$. Let $Y_1$ and $Y_2$ be two of the others,
and let $i\in\{1,2\}$. The set $Y_i$ contains $w-1$ of the cycles, and each cycle has at least
$D_*$ vertices. So $|Y_i|\ge(w-1)D_*\ge k'F(k')$. At most $2\cdot2s=k'$ edges leave $Y_i$, and
every component of $N[Y_i]$ has such an edge, because $b\notin Y_i$. So some component has a
vertex set $Y_i'$ with at least $F(k')$ vertices. The vertices of $Y_i$ have degree three in $N$,
so $N[Y_i']$ has at most $k'$ boundary vertices. By induction, every even set of boundary
vertices of $N[Y_i']$ has forest proportion less than $\frac35$. Lemma~\ref{lem:product}, for
$N$, $T$, $P=\emptyset$ and the disjoint sets $Y_1'$ and $Y_2'$, shows that the forest
proportion of $T$ is at most $\frac9{25}$.

\step{The size of $F$.} We have $2^{s}=O(k^{2}\log k)$ and $s=O(\log k)$. So $n_1=O(k^{4}\log^{3}k)$
and $n_2=O\bigl(k^{3}\log k\cdot F(4s_k)\bigr)$, and the definition gives
$F(k)=O\bigl(k^{6}\log^{5}k+k^{5}\log^{3}k\cdot F(4s_k)\bigr)$. Hence $F(k)\le C_1k^{7}$ for an
absolute constant $C_1$, by induction on $k$.
\end{proof}

\begin{corollary}[Two regions]\label{cor:two}
Let $G$ be a cubic multigraph, and let $P$ be a set of vertices such that $G-P$ has girth at least
$D_*$. Suppose that $V(G)\sm P$ contains two disjoint sets $Y_1,Y_2$ and that, for some
$k_1,k_2\ge1$, each $G[Y_i]$ is connected, $|\delta Y_i|\le k_i$ and $|Y_i|\ge C_1k_i^{7}$. Then
$q(G,P)\le\frac9{25}$.
\end{corollary}

\begin{proof}
The multigraph $G[Y_i]$ has at most $k_i$ boundary vertices and girth at least $D_*$. Use
Theorem~\ref{thm:pieces} and Lemma~\ref{lem:product} with $T=\emptyset$.
\end{proof}

\subsection{Section 8 in the second form}\label{sec:secondform}

Let $D_*$ and $C_1$ be as in Theorem~\ref{thm:pieces}. The statements of
Section~\ref{sec:first} change as follows. In each of them, $G-P$ has girth at least $D_*$.
\begin{itemize}[leftmargin=1.2em]
\item \emph{A long chain.} For an integer $\Lambda\ge1$: if
$T>(\Lambda+2)\bigl(\lceil2^{8}C_1\Lambda^{8}/m\rceil+1\bigr)$, then $q(G,P)\le\frac9{25}$. In
the proof, put $w=\lceil2^{8}C_1\Lambda^{8}/m\rceil+1$. The only sets $Y_l$ to discard are the at
most $\Lambda$ sets that contain a vertex with a neighbour in $A$. A set $Y$ that remains has
$|\delta Y|\le2\Lambda$ and $|Y|\ge(w-1)m\ge2\Lambda\cdot C_1(2\Lambda)^{7}$. So a component of
$G[Y]$ has at least $C_1(2\Lambda)^{7}$ vertices, and Corollary~\ref{cor:two} applies with
$k_1=k_2=2\Lambda$.
\item \emph{Heavy families.} The bound is
$\mu\le\tfrac32K+3\Lambda\kappa\,(\Lambda+3)\bigl(\lceil2^{8}C_1\Lambda^{8}/m\rceil+1\bigr)$.
The proof does not change.
\item \emph{Packing.} For every integer $d\ge D_*$, the multigraph $G-P$ has at most
$2^{10}C_1\,2^{d}d^{10}$ pairwise disjoint cycles of length at most $d$. The proof is the second
step of the proof of Proposition~\ref{prop:packing}, and there is no induction: with
$\Lambda=d$, $m=D_*$, $\kappa=1$ and $K=d+1$,
\[
  \mu\;\le\;\tfrac32(d+1)+3d\,(d+3)\bigl(\lceil2^{8}C_1d^{8}/D_*\rceil+1\bigr)\;\le\;2^{9}C_1\,d^{10},
  \qquad n\le2^{d+1}\mu .
\]
\item \emph{The proof of the theorem.} Now $|S|\le2^{10}C_1\,2^{D}D^{11}\le r/x^{2}$, and the last
display is
$\mu\le\tfrac32(1+\eta)+3L\eta\,(L+3)\bigl(\lceil2^{12}C_1L^{8}/x\rceil+1\bigr)=\tfrac32+o(1)$.
\end{itemize}
In this form the numbers $\nu_\Lambda$, the condition~\eqref{eq:dstar}, and part~(ii) of
Proposition~\ref{prop:packing} do not exist. The price is the longer proof of
Theorem~\ref{thm:pieces}, and the section is longer by 18 lines than it is now
(Section~\ref{sec:replaces}).

\section{The questions of the task}\label{sec:answers}

\subsection{Priority 1: one family of short cycles at one scale}

\emph{(1) Is there a direct bound for the number of vertices that we remove?} Yes.
Proposition~\ref{prop:packing} gives $|S|\le2^{D}D^{\,C_0+10}$ from $q(G,P)\ge\frac25$ and the
girth $D_*$, for one maximal family at the scale $D$. In the first proof, ``direct'' has one
restriction: the bound for the length $d$ uses the bound for the length $C_0\log(d+2)$. In the
second proof the packing bound has no induction.

\emph{(2) Do the pairs of removed short cycles with two links give the long chain?} Yes. This is
the proof of Proposition~\ref{prop:packing}: the family of the short cycles is a family for
Proposition~\ref{prop:heavy}. The proof of \M{Proposition~8.7} already has this idea. The new
point is that no deleted set is necessary for the girth of the blocks.

\emph{(3) Can we use the separators before the logarithmic girth is known?} Yes.
Proposition~\ref{prop:heavy} has no hypothesis on the girth of the family. The girth is in two
places only: in the number $\nu_\Lambda(G,P)$ of the blocks that we discard, and in
\M{Proposition~3.5}, which controls the second moment of the cycles of moderate length.

\emph{(4) Where does the present argument need two scales, and is this necessary?}
Section~\ref{sec:present} gives the place: Proposition~8.5 takes the girth of a block from a
deleted set, and the deleted set of the scale $x/16$ has more vertices than the chain has members.
It is not necessary. A block needs no short cycle \emph{inside} it, and the number of the blocks
with a short cycle is at most $\nu_L(G,P)=(\log r)^{O(1)}$.

\subsection{Priority 2: one structural lemma}

The task asks for a lemma of this form: either the forest probability is below $\frac25$, or the
graph contains two large disjoint regions with small boundaries. Proposition~\ref{prop:heavy}
with Corollary~\ref{cor:two} is this lemma. Suppose that a family has the properties~(a), (b)
and~(c) and that its weight is above the bound. Then
\begin{itemize}[leftmargin=1.2em]
\item the pairs with two links have weight $W_2<\frac23\mu^{2}-K\mu$, and
$q(G,P)<\frac25$ by the second moment; or
\item $V(G)\sm P$ contains two disjoint connected sets with at most $2\Lambda$ boundary edges and
at least $C_1(2\Lambda)^{7}$ vertices, and $q(G,P)\le\frac9{25}$ by Corollary~\ref{cor:two}.
\end{itemize}
The regions come from the pairs with two links, through \M{Lemma~8.6} and
Lemma~\ref{lem:blocks}. No packing bound is used.

The parts of this lemma are not all new. The second moment, \M{Lemma~8.6} and the blocks are
in~[M]. Three points are new: the statement for a general weighted family, which serves the short
cycles and the cycles of moderate length; the count of the blocks with a short cycle in the place
of a deleted set (first proof); and Theorem~\ref{thm:pieces}, which removes the condition on the
girth of the regions (second proof).

\subsection{Priority 3: the hypotheses of the final application}\label{sec:p3}

In Case~5 of \M{Section~9}, $G$ has no cyclic cut of size at most four, its girth is at least
$D_*$, and $P$ is a connected set of eight vertices. The task asks if these hypotheses exclude the
pairs with two links. They do not.

\begin{lemma}[Necklaces]\label{lem:necklace}
Let $R$ be a bipartite cubic graph with $n\ge10$ vertices and a Hamiltonian cycle
$v_0v_1\cdots v_{n-1}$, and let $M\ge4$. Let $G_M$ be the graph on the vertices $(v_a,i)$ with
$i\in\mathbb Z_M$ that has the edges $(v_a,i)(v_{a+1},i)$ and, for each chord $v_av_b$ of the
Hamiltonian cycle with $a$ even, the edges $(v_a,i)(v_b,i+1)$. The \emph{ring} $Z_i$ is the cycle
on the vertices $(v_a,i)$. Then
\begin{enumerate}
\item[\textup{(i)}] $G_M$ is a connected cubic graph, and its girth is at least the girth of $R$;
\item[\textup{(ii)}] $G_M$ has no cyclic cut of size at most four;
\item[\textup{(iii)}] two rings that are not consecutive have exactly two links.
\end{enumerate}
\end{lemma}

\begin{proof}
(i) The map $(v_a,i)\mapsto v_a$ is a covering map from $G_M$ to $R$. So a cycle of $G_M$ maps to
a closed walk of $R$ without an immediate return, and such a walk has at least as many edges as a
shortest cycle of $R$.

(ii) Each vertex is on one edge outside its ring. For a vertex $(v_a,i)$ with $a$ even this edge
goes to $Z_{i+1}$, and for $a$ odd it goes to $Z_{i-1}$. So $n/2\ge5$ edges join two consecutive
rings, and a ring is an induced cycle. Let $X$ be a set of vertices with $|\delta X|\le4$ such
that $X$ and its complement are not empty. We show that $G[X]$ or $G-X$ has no cycle. Call a ring
\emph{split} if it has vertices on both sides. A cut contains an even number of edges of a cycle.
So $\delta X$ contains at least two edges of a split ring, and at most two rings are split. Two
consecutive rings that are not split are on the same side.

If no ring is split, all rings are on one side. This is not possible. If one ring is split, the
other rings are consecutive, so they are on one side. The other side is a proper subset of a
ring, and it induces no cycle. Let two rings be split. Then $\delta X$ contains two edges of
each of them and no edge between two rings. Suppose that the split rings are not consecutive,
and let $Z_j$ be one of them. Its vertices with $a$ even are on the side of $Z_{j+1}$, and its
vertices with $a$ odd are on the side of $Z_{j-1}$. If these sides are equal, $Z_j$ is not split.
If they differ, all edges of $Z_j$ are in $\delta X$. Both cases are not possible. So the split
rings are $Z_j$ and $Z_{j+1}$. The other rings are consecutive and on one side. Let $Y$ be the
other side. A vertex of $Z_j$ with $a$ odd has an edge to $Z_{j-1}$, and this edge is not in $\delta X$.
So the vertices of $Y$ in $Z_j$ have $a$ even. They are consecutive on $Z_j$, because $\delta X$
contains two edges of $Z_j$. So there is one of them. In the same way, with the edges to
$Z_{j+2}$, the set $Y$ has one vertex in $Z_{j+1}$. So $|Y|=2$, and $Y$ induces no cycle.

(iii) Without two rings that are not consecutive, the graph has two components. Each of them is
joined to the two rings, and no edge joins the two rings.
\end{proof}

The Heawood graph ($n=14$, girth $6$) and the Tutte--Coxeter graph ($n=30$, girth $8$) are
bipartite, cubic and Hamiltonian. For every $g$ there is such a graph $R$ of girth at least $g$.
For example, take a cover of $R$ with the group $\mathbb Z_k$ in which one edge of the Hamiltonian
cycle has the voltage $1$ and the chords have the voltages $(2g)^{1},(2g)^{2},\dots$, for a large
$k$. The cover is bipartite and cubic, and the Hamiltonian cycle lifts to a Hamiltonian cycle. A
cycle of the cover with fewer than $g$ edges maps to a closed walk of $R$ whose voltages add up
to zero. For a large $k$ this is possible only if the walk uses each of these edges the same
number of times in the two directions. A shortest cycle of $R$ is not such a walk,
because it uses a chord one time, or it is the Hamiltonian cycle. So the girth increases, and we
repeat the step.

So there are cubic graphs of every girth, without a cyclic cut of size at most four, that contain
$M$ pairwise disjoint cycles of one length $n$ in which all pairs that are not consecutive have
exactly two links. For the family of these $M$ rings, with equal weights, $W_2=(1-3/M)\mu^{2}$.
The second moment of this family gives no bound below $\half$. For $M>2^{16n}$ the rings are
cycles of length less than $(\log r)/16$, so they are short cycles at the scale of the proof. The
theorem holds for these graphs. But a proof must find the two pieces between the rings, and the
hypotheses do not remove this case. The reason is that the hypothesis is on cuts of size at most
four, and the regions of the proof have at least $D_*\ge7$ boundary edges.

Two more remarks on the final application.
\begin{itemize}[leftmargin=1.2em]
\item Case~5 needs less than $q(G,P)<\frac25$. For a connected set $P$ of $p$ vertices the
argument of \M{Section~9} gives $c(G)\le\bigl(q(G,P)+2^{-1-p/2}\bigr)2^{r}$. So a bound
$q(G,P)\le\half-\varepsilon$, for one fixed $\varepsilon>0$ and connected sets of a fixed size, is
sufficient. I found no simpler proof for a weaker constant: one piece gives no bound below
$\half$, and two pieces give $\frac9{25}$.
\item The proofs of this note use the set $P$ in two ways only: it contains the vertex $p$, and
the regions avoid it. The size of $P$ is used in \M{Proposition~3.5} only.
\end{itemize}

\section{What stays, and why}\label{sec:stays}

\paragraph{The packing bound stays.} Proposition~8.7 of~[M] is not removed; it has a shorter
statement and no deleted set. It cannot be removed in this method. The second moment of
\M{Proposition~3.5} needs a graph $J$ of girth about $(\log r)/16$: the number of the links of two
cycles of length about $\log r$ is controlled by the girth (Step~2 of that proof). And the weight
of the cycles of $J$ is large only if few edges join $J$ to the deleted set (\M{Lemma~3.4} needs
$L\,d_J/v=o(1)$). So the case of many disjoint short cycles is a separate case, and the packing
bound is this case.

\paragraph{An induction over scales stays.} There are three facts.
\begin{enumerate}[leftmargin=1.6em]
\item The second moment alone gives the bound $\half$ and no more, because two regions with two
links have a correlation of two. A bound below $\half$ needs two pieces.
\item A piece needs a large layer. In the proof of \M{Lemma~8.2} the layer comes from a ball
that is a tree, of radius about $2\log k$ for $k$ boundary vertices. So a block with a boundary
of size $k$ must be free of cycles of length about $4\log k$, or we must count such blocks.
\item \M{Theorem~7.1} treats the graphs with a cycle shorter than a fixed number. Its error terms
come from the forest estimate with $\varepsilon=2^{-g-2}$. So it cannot treat a girth that grows
with $r$.
\end{enumerate}
The girth of the hypothesis is the constant $D_*$, and the final step needs blocks without cycles
of length about $C_0\log\log r$. An argument must pass from the one to the other. In~[M] it is the
sequence of the scales $2^{j}D_*$, with about $\log\log r$ steps and nested deleted sets. In the
first proof it is the induction $d\mapsto C_0\log(d+2)$ in Proposition~\ref{prop:packing}, with
about $\log^{*}r$ steps and no deleted set. In the second proof it is the induction
$k\mapsto4s_k$ in Theorem~\ref{thm:pieces}. I found no proof without such an induction.

\paragraph{The factor two is a parity.} The script \texttt{checks/necklace\_parity.py} examines
a necklace of $M$ rings of length four, for $4\le M\le7$, by complete enumeration. Let the
\emph{parity} of $\xi$ be the parity of its number of edges between two fixed consecutive rings.
The parity is even with probability $\half$. If it is odd, no ring is a component of $\xi$. If it
is even, each ring is a component with two times its probability without the condition, and for
rings that are pairwise not consecutive the probability that none of them is a component is the
product, as for independent events. This is a test on one family of graphs and not a theorem. It shows why the second moment stops at $\half$, and why the blocks of
Theorem~\ref{thm:pieces} need another argument when the parity is odd: there the cycles of the
chain are not free.

\paragraph{The lemma on pieces is false for girth four.} In a ladder with $n$ rungs, let $T$
consist of the two ends of one rail. Then every $T$-join is a path, for every $n$. So the forest
proportion is $1$, and the ladder has four boundary vertices and any number of vertices
(\texttt{checks/tjoin\_check.py} tests $n\le10$). In the proof of Theorem~\ref{thm:pieces} the
ladder is the case of cycles that are not free, and the induction does not start: the cycles of
length four have four boundary edges, and the pieces between them have four boundary vertices
again. The proof of Theorem~\ref{thm:pieces} needs $D_*=2s_{k_0}+1$, which is more than $200$.

\paragraph{An open question, and what it would remove.} I do not know the least girth for which
Theorem~\ref{thm:pieces} holds. In its proof the girth is used in the base only, for pieces with
fewer than $k_0$ boundary vertices. So the theorem holds at a smaller girth if it holds there for
these pieces, with some bound for the number of vertices. The pieces of the application are
induced subgraphs of a cubic multigraph without a cyclic cut of size at most four, and such a
multigraph has girth at least five. Suppose that the base holds for these pieces at girth five.
Then Theorem~\ref{thm:half} holds with $D_*=5$ in the final application, Case~4 of \M{Section~9}
is empty, and \M{Section~7} with \M{Lemma~5.5} and \M{Corollary~5.6} is not necessary. For two
boundary vertices, \M{Corollary~5.6} is a statement of this kind, with $\half+\varepsilon$ in the
place of $\frac35$. For more boundary vertices I have no proof and no counterexample. This is a
question and not a result.

\paragraph{Approaches that failed.}
\begin{itemize}[leftmargin=1.2em]
\item \emph{A marker without layers.} If a small connected set contains all boundary vertices of
a block, then \M{Theorem~3.1} gives a forest proportion of at most $\half+o(1)$ at once, with no
hypothesis on the girth. For a block between two cycles $C_i$ and $C_j$ the set is the two cycles
and a path between them. The obstruction is a long thin block: the path has almost as many
vertices as the block. The ladder is the extreme case.
\item \emph{A second moment in a block.} The cycles of the chain that lie in one block are
pairwise disjoint, so their total weight is at most $\eta$. They give no bound.
\item \emph{A chain that is longer than a trivial count.} The proof gives a chain of about
$2^{s}/L^{3}$ members, and this is less than $r^{1/16}$ for a girth of at most $x/8$. Without a
packing bound, the number of the disjoint short cycles has only a bound of order $r$. So we
cannot count the blocks with a short cycle without a packing bound.
\item \emph{A higher girth in \M{Theorem~7.1}.} See fact~3 above.
\end{itemize}

\section{Checks and files}\label{sec:checks}

All files are in \texttt{paper/minimal/}.
\begin{center}
\small
\begin{tabular}{@{}p{5.6cm}p{8.6cm}@{}}
\toprule
File & Content and result\\
\midrule
\texttt{research/apply\_one\_scale.py} & Applies the first proof to a copy of the paper
(Sections~8 and~1). Each replacement must match one time. The copy builds: 42 pages, 1615
lines, no undefined reference, no overfull line.\\\addlinespace
\texttt{research/apply\_two\_regions.py} & Applies the second proof to a copy. The copy builds:
42 pages, 1644 lines, no undefined reference, no overfull line.\\\addlinespace
\texttt{research/trial-one-scale.pdf}\newline \texttt{research/trial-fixed-girth.pdf} &
The two copies of the paper. Section~8 starts on page~34.\\\addlinespace
\texttt{checks/tjoin\_check.py} & Lemma~\ref{lem:free}\,(i) and the probabilities of
Lemma~\ref{lem:free} for one and two regions, on 300 random subcubic graphs with 8 to 13
vertices; the ladders with 3 to 10 rungs. Exit code~0 for the seeds 1 and 7.\\\addlinespace
\texttt{checks/tube\_cyclic\_cuts.py} & Lemma~\ref{lem:necklace} for the Heawood graph and
$M=4,5$: girth~6, every set of at most four edges tested, two links between the rings~0 and~2.
Exit code~0.\\\addlinespace
\texttt{checks/necklace\_parity.py} & The parity statement of Section~\ref{sec:stays} for
$4\le M\le7$. Exit code~0.\\
\bottomrule
\end{tabular}
\end{center}
The scripts test the statements on small graphs. They do not test the proofs of
Sections~\ref{sec:first} and~\ref{sec:second}, whose constants are too large for a computer.

\paragraph{The choice.} The first proof is the smaller change: \M{Lemma~8.2} stays, and the
section is shorter. The second proof has the cleaner statements and the stronger lemma, and the
section is longer. Each of them needs a decision by KNT before the manuscript changes, and each
makes the text differ more from the formal proof. My recommendation is the first proof for the
manuscript. Theorem~\ref{thm:pieces} then stays in this note; it is the start for the open
question of Section~\ref{sec:stays}.

\begin{thebibliography}{9}\small
\bibitem[M]{M} KNT, Astra and Fable~5.2, \emph{The maximum number of cycles in a graph of large
rank}, draft of 12 November 2026, 42 pages, file
\texttt{paper/minimal/exact-maximum/exact-maximum.pdf}.
\end{thebibliography}

\end{document}
